Several limitations of our approach warrant discussion. First, our cosmic-ray rejection model was trained exclusively on HST/WFC3 imaging, and its performance on ground-based data with different detector characteristics remains untested. Second, the false-positive rate of 2.3\% — while low — is dominated by compact, high-surface-brightness sources such as unresolved stars, which could bias source catalogs in crowded fields. Third, we have not explored the impact of varying exposure times on model performance. These limitations notwithstanding, our results demonstrate the viability of learned cosmic-ray rejection. Future work will delve into domain adaptation techniques to extend the model to Roman and Euclid imaging, and will investigate uncertainty quantification for the rejection masks. Ultimately, we envision a unified framework for artifact detection across all major imaging surveys.
